\documentclass[10pt]{article}
\usepackage[cmex10]{amsmath}
\usepackage{amssymb,amsthm}
\DeclareFontShape{OMX}{cmex}{m}{n}{<->cmex10}{}
\DeclareMathSizes{10}{10}{7}{7}
\DeclareMathSizes{9}{10}{7}{7}
\usepackage[margin=0.9in]{geometry}
\usepackage[hidelinks]{hyperref}
\newtheorem{theorem}{Theorem}
\title{A smallest fractional separation for spectral Chv\'atal}
\author{}
\date{November 2026}
\begin{document}
\maketitle

\begin{abstract}
For a downset on at most six elements, the Lov\'asz theta number of its
disjointness graph equals the size of its largest star. All these downsets
admit a clique partition of that size except one, up to relabeling. The
exception is the face family of the six-vertex triangulation of the real
projective plane. Its largest star has size $11$, its fractional independence
number is $35/3$, and its minimum clique partition has size $12$. The theta
equality also holds for every rank-three extension of this family by one element.
\end{abstract}

Let $G(\mathcal D)$ be the simple graph on the nonempty members of a finite
downset $\mathcal D$, with adjacency given by disjointness. Write $s(\mathcal D)$
for the largest number of sets through one element, with $s=\vartheta=0$
for empty support. Ellis, Filmus, and
Friedgut's Conjecture H is equivalent to $\vartheta(G(\mathcal D))=s(\mathcal D)$
\cite{eff}. They give a seven-element counterexample to the stronger
fractional independence bound.

\begin{theorem}
If $\mathcal D$ has support of size at most six, then
$\vartheta(G(\mathcal D))=s(\mathcal D)$. A clique partition of size
$s(\mathcal D)$ exists except when $\mathcal D$ is the closure of the ten
triples in \eqref{eq:facets}, up to relabeling. This exception has
$s(\mathcal D)=11$, fractional independence $35/3$, and clique partition
number $12$. In particular, six is the least support size for a fractional
gap, and the largest fractional gap through support size six is $2/3$.
\end{theorem}

\section{Reductions}
\paragraph{Singleton deletion.}
In a graph $G$, let $v$ and all its nonneighbors form an independent set
$S$ of size $k$. Then
\[
\vartheta(G)=\max\{k,\vartheta(G-v)\}.
\]
The lower bound is monotonicity and independence. For the upper bound take
$t\geq\max\{k,\vartheta(G-v)\}$. The standard theta dual gives a Gram matrix
on $G-v$ with diagonal $t-1$ and off-diagonal entries $-1$ on nonedges.
Write $w_u$ for its Gram vectors. Give $v$ the vector
\[
-\frac{1}{t-k+1}\sum_{u\in S\setminus\{v\}}w_u.
\]
Its inner product with each $w_u$ in the sum is $-1$. Its squared norm is
$(k-1)/(t-k+1)\leq t-1$, so an orthogonal coordinate supplies the missing
squared norm. This gives the required dual matrix on $G$. Each singleton
in $G(\mathcal D)$ satisfies the hypothesis because its nonneighbors all
contain its element. Apply the identity one singleton at a time. Thus
Conjecture H holds for $\mathcal D$ exactly when the induced graph on its
sets of size at least two has theta at most $s(\mathcal D)$.

If the graph on the larger sets has a partition into at most $s$ cliques,
pad the partition to $s$ slots. At most $s-1$ larger sets contain any one
element, so its singleton fits in a slot containing no such set. Each pair
of singletons is disjoint, so these choices can all be made independently.

\paragraph{Complementary pairs.}
Let $U$ be the support. If $U$ belongs to the downset, it is the full cube.
Pair its proper nonempty sets with their complements and put $U$ alone in
one clique. This gives $2^{|U|-1}$ cliques when $U$ is nonempty. Otherwise
remove every complementary pair of sets present in $\mathcal D$. The
remaining family is a downset. Indeed, if $A$ remains and $B\subseteq A$
had been removed, then $U\setminus B$ was present and its subset
$U\setminus A$ would have been present too. Each removed pair reduces every
star by one, and its two sets form a clique. Thus a clique partition lifts
by appending the removed pairs. A theta upper bound lifts by adding at most
one per pair, by subadditivity over a vertex partition \cite{lovasz}.
Repeat if the support decreases. Consequently it suffices to consider
families with no complementary pair on their actual support.

\section{Exhaustive check}
On a fixed five-element set, the Python enumerator generates each antichain
of nonempty facets once, giving $7\,579$ nontrivial downsets including all
smaller supports. After the two reductions the graph has clique size at
most two. A maximum matching gives a minimum clique partition. The program
recurses on one vertex, trying a singleton or every available disjoint
partner. Every resulting bound is at most the remaining star size.

On six elements the C++ enumerator generates antichains of proper facets
with full support. It excludes a pair of facets whose union is the whole
support. This condition is equivalent to the absence of complementary
pairs in their closure. It yields $1\,379\,272$ labeled downsets. In the
graph of sets of size at least two, every nontrivial clique includes a
two-set, because disjoint triples would be complementary. Cliques of size
three consist of three pairs. A recursive cover search selects one pair
and tries its singleton block, all disjoint partners, and all three-pair
blocks containing it. A block of size $q$ saves $q-1$ slots. A branch can
save at most the number of remaining pairs and at most two thirds of the
number of remaining vertices. These bounds justify the two pruning tests.
Exactly twelve labeled families fail the target cover. Replaying their
orbit identifies them with the closure of
\begin{equation}\label{eq:facets}
123,\ 124,\ 135,\ 245,\ 345,\ 236,\ 146,\ 346,\ 156,\ 256.
\end{equation}
The family contains every pair and one triple from each complementary
triple pair. An additional proper set on these six elements would have a
complement already in the family. Removing complementary pairs in a
strict superfamily therefore cannot leave this family intact. The
full-cube case has already been covered. Thus all cases outside this
orbit have an integral clique partition of the required size.

\section{The exceptional family}
The family in \eqref{eq:facets} has six singletons, fifteen pairs and ten
triples. Each atom belongs to five pairs and five triples, so $s=11$.
The link of each vertex is a five-cycle. The complex is a connected closed
surface with Euler characteristic $6-15+10=1$, giving the real projective plane.
On its $25$ nonsingletons write $A_{22}$ for the adjacency matrix of
disjoint pair--pair incidences and $A_{23}$ for the adjacency matrix of
disjoint pair--triple incidences, both embedded in the full matrix. The
theta dual at $t=11$ is certified by
\[
M=11I-J+2A_{22}+4A_{23}.
\]
Its characteristic polynomial, computed in exact arithmetic, is
\[
(x-11)^5(x^2-24x+127)^4(x^2-18x+1)^5(x^2-9x+2).
\]
It has no nonpositive real roots. Since $M$ is real symmetric, it is
positive definite. The accompanying checker also verifies positivity
directly by fraction-free elimination. Singleton deletion gives the full
theta upper bound of $11$, and a star attains it.

For fractional independence, give singleton, pair and triple vertices
weights $0,1/3,2/3$, respectively. A disjoint clique has at most three
pairs, or at most one triple and one pair. Thus these weights are feasible
and sum to $35/3$. For the matching upper certificate, assign weight
$1/3$ to each partition into a facet, a pair in its complement, and the
remaining singleton. Assign weight $1/9$ to each perfect matching of the
six elements into pairs. Every facet is covered to weight one. Each pair
is disjoint from two facets and belongs to three perfect matchings, so its
coverage is also one. Each singleton has coverage $5/3$. The total
covering weight is $30/3+15/9=35/3$. The repository verifies these weights
and an explicit twelve-clique partition, which is minimum by the
fractional lower bound. The full conjecture H remains open.

\newpage
\section{One-point extensions}
Suppose $\mathcal D'$ has rank at most three and its restriction to six
elements is \eqref{eq:facets}. The triples containing a seventh element
define a graph $L$ on the original six elements, with an edge for each
pair appearing in such a triple. The pairs containing the seventh element
must include every endpoint of $L$. Pairs to isolated vertices of $L$
are optional. When all six pairs are present, the code checks all
$688$ orbits of $L$ under the automorphisms of the base family. A rational
theta dual on the nonsingletons certifies each orbit.

If $L$ is nonempty, the largest old star has size $12+\Delta(L)$ regardless
of the optional pairs. The new star has size at most $7+|E(L)|$. If the
new star in the all-pairs extension is no larger, deleting optional pairs
leaves the star maximum unchanged, so the same theta bound follows by
monotonicity. If an optional pair exists and $7+|E(L)|>12+\Delta(L)$, then
$2|E(L)|\leq5\Delta(L)$ and $|E(L)|\geq6+\Delta(L)$ force $L$ to be $K_5$
with one isolated vertex. The code checks the one missing-pair case with
an additional rational dual. The base automorphism group is transitive on
the six elements, so the choice of isolate makes no difference. For an
empty $L$, retaining any optional pair keeps the star maximum at $12$.
Retaining none adds at most a singleton, which is covered by singleton
deletion. This proves the claimed rank-three extension result.

\section{Replaying the verification}
The supplementary repository includes the exhaustive enumerators, the
twelve undecided labeled forms, and the 689 dual certificates used in
Section 4. The replay commands in the README require only Python and a
C++17 compiler. The exact matrix checker uses integer entries and
fraction-free elimination. The optional numerical search proposes
candidate matrices and uses the same exact checker before accepting them.

\begin{thebibliography}{9}
\bibitem{eff} D. Ellis, Y. Filmus, and E. Friedgut.
Chv\'atal's conjecture: a proof from The Book. October 2026 revision.
\href{https://arxiv.org/html/2609.28404v3\#S5}{arXiv:2609.28404v3, Section 5}.
\bibitem{lovasz} L. Lov\'asz. On the Shannon capacity of a graph.
\emph{IEEE Transactions on Information Theory} 25 (1979), 1--7.
\href{https://doi.org/10.1109/TIT.1979.1055985}{doi:10.1109/TIT.1979.1055985}.
\end{thebibliography}
\end{document}
